\documentclass{report}
\usepackage[utf8]{inputenc}
\usepackage[T1]{fontenc}
\usepackage{geometry}
\usepackage{graphicx}
\usepackage{float}
\usepackage{hyperref}
\usepackage[english,indonesian]{babel}

\title{Optimasi Pemrosesan Template LaTeX Menggunakan Pendekatan Berbasis AI}
\author{Academic Support Agent}
\date{\today}

\begin{document}
\maketitle

\begin{abstract}

\end{abstract}

\chapter{PENDAHULUAN}


\chapter{TINJAUAN PUSTAKA}


% Original content (commented out):
% \documentclass[a4paper,12pt]{report}	
\usepackage[brazil, english]{babel}   
\usepackage[utf8]{inputenc} 
\usepackage{indentfirst}
\usepackage{mathptmx}

\input{TesePack}

\include{abbreviations}

\begin{document}

\thispagestyle{empty}

\input{Capa}

\pagebreak           

\pagenumbering{roman}
\setcounter{page}{2}

\input{Preambulo}

\input{modelo_latex/siglas}

\def\listfigurename{\begin{center}LISTA DE FIGURAS\end{center}}\listoffigures

\pagebreak\def\listtablename{\begin{center}LISTA DE TABELAS\end{center}}\listoftables

\pagebreak\def\contentsname{\begin{center}SUMÁRIO\end{center}}\tableofcontents

\cleardoublepage

\pagenumbering{arabic}

\input{cap1}

\input{cap2}

\input{modelo_latex/cap3}

\input{modelo_latex/cap4}

\input{modelo_latex/cap5}

\normalsize
\cleardoublepage
\addcontentsline{toc}{chapter}{Referências Bibliográficas}
\bibliographystyle{coppe}
\bibliography{biblio}

\end{document}


\end{document}